% Generated draft: requires mathematical and source review.
\documentclass[11pt]{article}
\usepackage[margin=25mm]{geometry}
\usepackage{fontspec}
\setmainfont{Latin Modern Roman}
\usepackage{amsmath,amssymb,mathtools,unicode-math}
\setmathfont{Latin Modern Math}
\usepackage{enumitem,needspace,xurl,hyperref}
\setlength{\parindent}{0pt}
\setlength{\parskip}{5pt}
\setlength{\emergencystretch}{3em}
\newcommand{\N}{\mathbb{N}}
\newcommand{\Z}{\mathbb{Z}}
\newcommand{\Q}{\mathbb{Q}}
\newcommand{\R}{\mathbb{R}}
\newcommand{\C}{\mathbb{C}}
\newcommand{\F}{\mathbb{F}}
\newcommand{\A}{\mathbb{A}}
\newcommand{\PP}{\mathbb{P}}
\newcommand{\E}{\mathbb{E}}
\newcommand{\eps}{\varepsilon}
\newcommand{\dd}{\,\mathrm{d}}
\begin{document}
\section*{UnsolvedMath ID: 700821}
\textbf{Source problem code:} catalogue entry
\subsection{Definitions and mathematical statement}
For a finite group \(G\), its integral group ring is
\[
\Z[G]=\left\{\sum_{g\in G}a_g g:a_g\in\Z\right\},
\]
with coefficientwise addition and multiplication induced by the group multiplication and distributivity. Its group of units is
\[
U(\Z[G])=\{u\in\Z[G]:\text{there exists }v\in\Z[G]\text{ with }uv=vu=1\}.
\]
For a subgroup \(H\leq U(\Z[G])\), the index \([U(\Z[G]):H]\) is the number of left cosets of \(H\). The subgroup is proper if \(H\neq U(\Z[G])\).

\textbf{Part (a).} The supplied source attributes to H. Bass an explicit construction, for a finite cyclic group \(C\), of a proper subgroup \(B_C\) of finite index in \(U(\Z[C])\). Determine the index
\[
[U(\Z[C]):B_C].
\]
Here \(B_C\) denotes the particular subgroup in the cited Bass construction. The supplied record does not give its generators or another definition identifying it, nor does it specify restrictions on the order of \(C\). Thus this is a source-dependent formulation: a self-contained index problem requires the precise construction and its applicable hypotheses from the cited reference.

\textbf{Part (b), attributed to A. A. Bovdi and marked archived solved in the supplied source.} For each finite abelian non-cyclic group \(A\), construct subgroups analogous to Bass' subgroup. In the terms of the preceding description, the requested objects are explicitly constructed subgroups \(H_A\) satisfying
\[
H_A<U(\Z[A]),\qquad [U(\Z[A]):H_A]<\infty.
\]
The record does not specify further properties required by the word ``analogous''; these must be checked against Bass' construction and the original problem.

The supplied archive reports constructions for finite \(p\)-groups with \(p\neq2\) by K. Hoechsmann and J. Ritter, and a general construction by R. Zh. Aleev for groups of central units of integral group rings of arbitrary finite groups, including nonabelian groups. Here the group of central units means
\[
U(\Z[G])\cap Z(\Z[G]),\qquad
Z(\Z[G])=\{z\in\Z[G]:zx=xz\text{ for every }x\in\Z[G]\}.
\]
These are supplied archive assessments, not independently verified mathematical conclusions.
\subsection{Short English statement}
Determine the index of the specific finite-index subgroup of the units of a finite cyclic integral group ring constructed by Bass. The source also records Bovdi's request for analogous constructions for finite abelian non-cyclic groups as archived solved. The precise Bass subgroup and the full meaning of ``analogous'' require the cited sources; the supplied registry labels the combined problem partially solved.
\subsection{Sources}
\begin{itemize}
\item Supplied local UnsolvedMath registry record: permanent ID \(700821\), source code \texttt{catalogue entry}, title \emph{Kourovka Notebook Problem 12.1}; upstream revision \texttt{d371ee917ac9647fbb06cd6fc529deec6f3dda82}. Proposed by R. Zh. Aleev in 1992; part (b) is attributed to A. A. Bovdi.
\item \emph{The Kourovka Notebook (Unsolved Problems in Group Theory)}, 12th issue (1992), Problem 12.1; the record describes an update in the 21st issue (2026), \texttt{arXiv:1401.0300v48}. The supplied collection description names E. I. Khukhro and V. D. Mazurov as editors.
\item H. Bass, \emph{Topology} \textbf{4}, no. 4 (1966), 391--400. Cited in the source for the explicit subgroup construction. No article title or defining construction is supplied.
\item K. Hoechsmann and J. Ritter, \emph{J. Pure Appl. Algebra} \textbf{68}, no. 3 (1990), 325--333. The supplied Archive of Solved Problems credits this reference with constructions for \(p\)-groups with \(p\neq2\).
\item R. Zh. Aleev, \emph{Mat. Tr. Novosibirsk Inst. Math.} \textbf{3}, no. 1 (2000), 3--37 (Russian). The supplied archive credits this reference with the general case of groups of central units of integral group rings of arbitrary finite groups, not necessarily abelian.
\item The registry status review is dated \texttt{2026-10-16} and gives the canonical status as \texttt{partially\_solved}, with evidence marked ``published'' and a note that solved parts occur in the archive while remaining parts occur in the main section. This is supplied dated curation, not independently verified source or literature assessment. The background supplies the fuller archive account; the status-review quotation is truncated.
\end{itemize}
\end{document}
